\section{Security Properties}
\label{sec:claims}

We state six security properties of cMCP. Each property corresponds to a reproducible
experiment in Section~\ref{sec:eval}. The properties are stated as informal invariants
over the cMCP runtime; formal proofs are left for future work.

\begin{property}[TEE-Measured Policy Enforcement]
\label{prop:policy-hash}
Let $B$ be the Cedar policy bundle loaded at TEE startup, and let $h_B$ be its
SHA-256 hash recorded in TEE-protected memory. For any tool call processed after startup:
\begin{enumerate}[leftmargin=1.5em,label=(\alph*)]
  \item The gateway starts only if the hash of the loaded bundle equals the configured
    expected value; a bundle swapped on disk before startup raises
    \texttt{PolicyHashMismatch}. The measured revision does not re-hash the bundle in
    memory before each call.
  \item The TRACE Claim for the session carries $h_B$ in the \texttt{policy.bundle\_hash}
    field. A verifier who independently hashes the policy files can confirm they match.
  \item Altering the bundle hash in a signed TRACE Claim invalidates its Ed25519
    signature. A claim that the gateway itself signs over a different bundle verifies; it
    is caught by comparing its \texttt{policy.bundle\_hash} with the approved value.
\end{enumerate}
\end{property}

\begin{property}[Session Sensitivity Tracking]
\label{prop:session-sensitivity}
Let $S$ be a session, and let $T_{\text{phi}}$ be the first tool call in $S$ whose
response is tagged with a sensitivity level in
$\{\text{hipaa\_phi}, \text{mnpi}, \text{trade\_secret}\}$.
\begin{enumerate}[leftmargin=1.5em,label=(\alph*)]
  \item After $T_{\text{phi}}$, \texttt{max\_sensitivity} is raised to the highest tag
    observed and never decreases below that level for the remainder of $S$.
  \item Any subsequent Cedar evaluation in $S$ that uses the current session sensitivity as
    context will observe the raised level, allowing policies to enforce cross-boundary
    restrictions on all calls after $T_{\text{phi}}$.
  \item Per-call sensitivity evaluation (evaluating each call independently without session
    context) cannot enforce such restrictions, because the PHI sensitivity is in the session
    state, not in the parameters of the subsequent call.
\end{enumerate}
The operational cost of this guarantee is a nonzero false-discovery proportion: calls after
$T_{\text{phi}}$ that are provably unrelated to the PHI are nonetheless flagged.
We quantify this cost empirically in Section~\ref{sec:eval}.
\end{property}

\begin{property}[Tool Catalog Rug-Pull Detection]
\label{prop:rug-pull}
Let $C$ be the approved tool catalog loaded and hash-pinned at TEE startup. For any tool
call to tool $t$:
\begin{enumerate}[leftmargin=1.5em,label=(\alph*)]
  \item If $C$ still hashes to its startup value and $t$ is in $C$, the call proceeds to
    Cedar evaluation, and the agent only ever sees $t$'s approved definition.
  \item If $C$ differs from the approved hash at load, \texttt{CatalogHashMismatch} is
    raised and the gateway does not start. If the in-memory catalog changes after startup,
    the call is refused before dispatch and is not propagated to the upstream server.
  \item The per-tool hash function has the avalanche property: a single-character change
    in any field changes approximately half the bits of the output hash.
\end{enumerate}
\end{property}

\begin{property}[TRACE Claim Key Binding]
\label{prop:nonce}
Let $K$ be the TEE keypair generated at startup. The attestation nonce is
$n = \text{TP}(K_{\text{pub}}) \,\|\, s$, where $\text{TP}(K_{\text{pub}})$ is the
32-byte RFC~7638 JWK Thumbprint of the public key and $s$ is a 32-byte random salt
drawn once per enclave startup. The report a session claim carries uses
$s = \text{SHA-256}(\text{chain\_root})$ instead. The nonce is committed into the hardware
report's \texttt{report\_data} field. The session identifier $\sigma$ is also covered by
the Ed25519-signed claim body.
\begin{enumerate}[leftmargin=1.5em,label=(\alph*)]
  \item $\text{TP}(K_{\text{pub}})$ is deterministic and re-derivable by any verifier
    from the \texttt{cnf.jwk} public key, so the verifier confirms
    $\text{report\_data}[{:}32] = \text{TP}(K_{\text{pub}})$ (key binding).
  \item The salt $s$ makes nonces distinct across startups that share a key, with
    overwhelming probability. Freshness against a verifier's own challenge is not
    provided by $s$ and requires a verifier-supplied nonce.
  \item Distinct gateway instances produce distinct key thumbprints
    ($K_1 \neq K_2 \Rightarrow \text{TP}(K_{1,\text{pub}}) \neq \text{TP}(K_{2,\text{pub}})$
    with overwhelming probability), so a claim cannot be presented under a different
    TEE instance's identity.
  \item Session binding: $\sigma$ is covered by the Ed25519 signature over the claim body.
    Replacing $\sigma$ invalidates the signature, so a claim valid for session $\sigma_1$
    cannot be replayed as $\sigma_2$.
  \item Modifying any field of a TRACE Claim invalidates the Ed25519 signature computed
    over the canonical JSON serialization of the claim body.
  \item Removing audit entries from an exported, signed audit bundle changes its hash and
    invalidates the export signature. A verifier holding the claim can also recompute the
    chain tip from the entries it receives and compare it with the signed tip.
\end{enumerate}
\end{property}

\begin{property}[Temporal Adjacency Provenance Approximation]
\label{prop:temporal}
Let $L$ be the \texttt{SessionCallLog} for session $S$. Each entry $e_i \in L$ has a
monotonically increasing sequence number $\text{seq}(e_i)$.
\begin{enumerate}[leftmargin=1.5em,label=(\alph*)]
  \item Sequence numbers are strictly increasing: $i < j \Rightarrow \text{seq}(e_i) < \text{seq}(e_j)$.
  \item A cross-boundary event is recorded whenever a call to a tool in domain $D_2$ follows
    any call to a tool in a high-sensitivity domain $D_1 \in \{\text{phi}, \text{pii},
    \text{pci}, \text{restricted}\}$, and $D_2 \neq D_1$. (Current releases add
    \texttt{hipaa\_phi}, \texttt{mnpi} and \texttt{pci\_data} to this set.)
  \item Every call graph summary includes a \texttt{edges\_represent} field stating
    explicitly that edges represent temporal adjacency, not data provenance.
  \item Zero false negatives by construction: if call $B$ follows call $A$ and $A$ touched
    a high-sensitivity domain, then $\text{seq}(B) > \text{seq}(A)$, so $B$ is necessarily
    included in the set of calls adjacent to $A$.
  \item False positives are possible: a call $B$ that is temporally adjacent to a PHI call
    $A$ but provably carries no PHI in its parameters or response is still attributed as
    PHI-adjacent.
  \item Denied calls are recorded in the graph with their policy decision, because the
    agent's attempt to make the call is itself evidence of awareness of the PHI context.
\end{enumerate}
\end{property}

\begin{property}[Cross-Organizational Attestation Chain]
\label{prop:cross-org}
Let $(K_1, \sigma_1)$ be the gateway TEE keypair and session identifier (Phase~1), and
$(K_2, \sigma_2)$ be the server TEE keypair and session identifier (Phase~2), with
$\sigma_1 = \sigma_2 = \sigma$.
\begin{enumerate}[leftmargin=1.5em,label=(\alph*)]
  \item $K_1$ and $K_2$ are generated independently; they are distinct with overwhelming
    probability.
  \item Both claims reference $\sigma$, providing the linkage without sharing private key
    material.
  \item Each phase's nonce binds its own key: $n_1 = \text{TP}(K_{1,\text{pub}}) \| s_1$
    and $n_2 = \text{TP}(K_{2,\text{pub}}) \| s_2$ are distinct because $K_1 \neq K_2$.
    The shared $\sigma$, carried in both signed claim bodies, provides the linkage. (The
    software experiment uses the older form $\text{SHA-256}(K_{\text{pub}} \| \sigma)$; see
    Section~\ref{sec:eval}.)
  \item Each claim is independently verifiable against its own public key.
  \item Tampering with the Phase~1 claim invalidates its Ed25519 signature but does not
    affect the validity of the Phase~2 claim.
  \item A server binary swap (loading a different TEE image) changes the TEE measurement
    in the Phase~2 attestation report, which a verifier who checks the expected measurement
    will detect.
\end{enumerate}
\end{property}
